\documentclass[11pt]{article}
\usepackage[margin=1in]{geometry}
\usepackage{amsmath,amssymb,amsthm,xcolor,hyperref}
\hypersetup{colorlinks=true,allcolors=blue!45!black}
\newtheorem{theorem}{Theorem}
\newtheorem{corollary}[theorem]{Corollary}
\title{The Scaled Tepper Identity: An Affine Reduction and a Lean Certificate}
\author{Myeongjin Kang}
\date{October 25, 2026}
\begin{document}
\maketitle
\begin{abstract}
We give a complete proof of the scaled Tepper identity posed as Conjecture 4.5 by Bayat and Teimoori Faal. The positive-integer-step statement follows immediately from their unit-step theorem by changing variables. We also supply an independently checked Lean proof, based on finite differences, for every nonzero real step. This note claims no new-to-literature priority.
\end{abstract}
\section{Statement and proof}
\begin{theorem}\label{thm:main}
Let $P\in\mathbb R[t]$ have degree $n$ and leading coefficient $a_n$. For all $x\in\mathbb R$ and all $l\in\mathbb R\setminus\{0\}$,
\[
\sum_{k=0}^n(-1)^k\binom nkP(x-lk)=a_n l^n n!.
\]
\end{theorem}
\begin{proof}
Set $R(t)=P(lt)$. Since $l\ne0$, this polynomial has degree $n$ and leading coefficient $a_nl^n$. Theorem 4.3 of \cite{BT} applied to $R$ at $x/l$ gives
\[
\sum_{k=0}^n(-1)^k\binom nkR(x/l-k)=a_nl^nn!.
\]
The evaluation in each summand is $P(x-lk)$, giving the assertion.
\end{proof}
\begin{corollary}
Conjecture 4.5, equation (13), of \cite{BT}, as recorded in MathDB 354537, holds for every positive integer $l$.
\end{corollary}
\section{Finite-difference verification}
Here is a proof route that does not assume Theorem 4.3. For a function $F:\mathbb R\to\mathbb R$, define $\Delta F(t)=F(t+1)-F(t)$. Induction on $n$, using Pascal's identity, yields
\[
\Delta^nF(0)=\sum_{k=0}^n(-1)^{n-k}\binom nk F(k).
\]
For a polynomial $Q$ of degree $n$, taking a forward difference lowers its degree by one and multiplies the leading coefficient by its degree. Consequently $\Delta^nQ=n!\operatorname{lc}(Q)$, as a constant function. These are established results in the pinned Lean library, under the names \texttt{fwdDiff\_iter\_eq\_sum\_shift} and \texttt{Polynomial.fwdDiff\_iter\_degree\_eq\_factorial}.

For $Q(t)=P(x-lt)$, the degree is $n$ and the leading coefficient is $a_n(-l)^n$. Multiplying its forward-difference expansion by $(-1)^n$ proves Theorem~\ref{thm:main}, because for $k\le n$,
\[
(-1)^n(-1)^{n-k}=(-1)^k,\qquad (-1)^n(-l)^n=l^n.
\]
The checked code explicitly proves the sign identity and invokes the polynomial-composition degree and leading-coefficient theorems. The final declaration accepts an arbitrary polynomial, identifies its natural degree with $n$, and specializes $l$ to a positive natural number. It also handles the zero polynomial, for which both sides vanish.
\section{Status and reproducibility}
The archived MathDB record contains no posted solutions or formalizations. The retrieved original source labels the scaled formula as a conjecture; nevertheless it is an elementary consequence of the immediately preceding theorem. We therefore claim a complete answer and a local formalization, without asserting that this observation is absent from the literature.

The code is in \texttt{formal/Tepper.lean}; \texttt{formal/Checks.lean} inspects every declaration. The rebuild driver pins Lean 4.34.1 and mathlib revision \texttt{d13f23b7} (full hash in the supplied pins), deletes old local object files, compiles with warnings as errors, and verifies axiom reports. Only propositional extensionality, classical choice, and quotient soundness occur. No conjectural mathematical axiom or admitted proof is used.
\begin{thebibliography}{9}
\bibitem{BT} M. Bayat and H. Teimoori Faal, \emph{A Generalization of Tepper's Identity}, arXiv:2205.04237v1 (2022), Theorem 4.3 and Conjecture 4.5. \url{https://arxiv.org/html/2205.04237v1}.
\end{thebibliography}
\end{document}
